% Part: lambda-calculus
% Chapter: church-rosser
% Section: parallel-beta-reduction

\documentclass[../../../include/open-logic-section]{subfiles}

\begin{document}

\olfileid{lam}{cr}{pb}

\olsection{સમાંતર $\beta$-લઘુકરણ}

અહીં \emph{સમાંતર $\beta$-લઘુકરણ}નો ખ્યાલ રજૂ કરીને
તે ચર્ચ--રોસર ગુણધર્મ ધરાવે છે એમ સાબિત કરીશું.

\begin{defn}[સમાંતર $\beta$-લઘુકરણ, $\bredpar$] \ollabel{defn:bredpar}
  પદોનું સમાંતર લઘુકરણ ($\bredpar$) આગમનથી આ રીતે
  વ્યાખ્યાયિત થાય છે:
  \begin{enumerate}
    \item \ollabel{defn:bredpar1} $x \bredpar x$.
    \item \ollabel{defn:bredpar2} જો $N \bredpar N'$, તો
      $\lambd[x][N] \bredpar \lambd[x][N']$.
      \sourcecorrection{OLLAM-033}{સ્થિર અંગ્રેજી મૂળના
      બીજા નિયમની પૂર્વધારણા $N \xrightarrow{\beta} N'$
      છે, જ્યારે પછીની સ્વવાચકતા અને આગમનાત્મક
      સાબિતીઓ આ જ સંબંધની પૂર્વધારણા
      $N \bredpar N'$ વાપરે છે. અહીં તે સુધાર્યું છે.}
    \item \ollabel{defn:bredpar3} જો $P \bredpar P'$ અને
      $Q \bredpar Q'$, તો $PQ \bredpar P'Q'$.
    \item \ollabel{defn:bredpar4} જો $N \bredpar N'$ અને
      $Q \bredpar Q'$, તો
      $(\lambd[x][N])Q \bredpar \Subst{N'}{Q'}{x}$.
  \end{enumerate}
\end{defn}

સમાંતર $\beta$-લઘુકરણથી પદમાંનાં ગમે તેટલાં
સંકોચ્ય પદો એક જ પગલામાં સંકોચી શકાય છે. તે સામાન્ય
$\beta$-લઘુકરણથી જુદું છે: અહીં
મૂળ પદમાં પહેલેથી હાજર સંકોચ્ય પદો જ સંકોચી શકાય;
સમાંતર $\beta$-લઘુકરણથી નવાં ઊભાં થતાં સંકોચ્ય પદો એ જ
પગલામાં સંકોચી શકાતાં નથી. દાખલા તરીકે,
$(\lambd[f][fx])(\lambd[y][y])$નું એક પગલામાં
સમાંતર $\beta$-લઘુકરણ તેના પોતાના પદમાં અથવા
$(\lambd[y][y])x$માં થઈ શકે છે, પણ સીધું~$x$માં
નહીં. જોકે સામાન્ય $\beta$-લઘુકરણથી તે~$x$માં
જાય છે. નવું સંકોચ્ય પદ પહેલા સમાંતર $\beta$-લઘુકરણના
પગલા પછી જ ઊભું થાય છે; બીજા સમાંતર $\beta$-લઘુકરણના
પગલાથી~$x$ મળે છે.

\begin{thm}\ollabel{thm:refl}
  $M \bredpar M$.
\end{thm}

\begin{proof}
  અભ્યાસપ્રશ્ન.
\end{proof}

\begin{prob}
  \olref[lam][cr][pb]{thm:refl} સાબિત કરો.
\end{prob}

\begin{defn}[$\beta$-પૂર્ણ વિકાસ]\ollabel{defn:bcd}
  $M$નો \emph{$\beta$-પૂર્ણ વિકાસ}~$\bcd{M}$ આ રીતે
  આગમનથી વ્યાખ્યાયિત થાય છે:
  \begin{align}
    \bcd{x} &= x \ollabel{defn:bcd1} \\
    \bcd{(\lambd[x][N])} &= \lambd[x][\bcd{N}] \ollabel{defn:bcd2}\\
    \bcd{(PQ)} &= \bcd{P}\bcd{Q} && \text{જો $P$ એ $\lambd$-અમૂર્તીકરણ ન હોય}
    \ollabel{defn:bcd3} \\
    \bcd{((\lambd[x][N])Q)} &= \Subst{\bcd{N}}{\bcd{Q}}{x} \ollabel{defn:bcd4}
  \end{align}
\end{defn}

પદનો $\beta$-પૂર્ણ વિકાસ નામ મુજબ ``પૂર્ણ સમાંતર
લઘુકરણ'' છે. સામાન્ય સમાંતર $\beta$-લઘુકરણમાં કોઈ
સંકોચ્ય પદ ન સંકોચવાનું પસંદ કરી શકાય છે; પૂર્ણ વિકાસમાં
બધાંનું સંકોચન કરવું જ પડે. તેથી
$(\lambd[f][fx])(\lambd[y][y])$નો પૂર્ણ વિકાસ
$(\lambd[y][y])x$ છે, મૂળ પદ પોતે નહીં.

\begin{editorial}
  આ વ્યાખ્યામાં એક અડચણ છે: આપણે હજી
  ($\lambd$-)પદો પર વિધેયોનું પુનરાવર્તી વ્યાખ્યાયન
  કેવી રીતે કરવું તે રજૂ કર્યું નથી. આગળ સુધારીશું.
\end{editorial}

\begin{lem}\ollabel{lem:comp}
  $M \bredpar M'$ અને $R \bredpar R'$ હોય તો
  $\Subst{M}{R}{y} \bredpar
  \Subst{M'}{R'}{y}$.
\end{lem}

\begin{proof}
  $M \bredpar M'$ની !!{derivation} પર આગમનથી.
  \begin{enumerate}
    \item છેલ્લો નિયમ \olref{defn:bredpar1} હોય: અભ્યાસપ્રશ્ન.
    \item છેલ્લો નિયમ \olref{defn:bredpar2} હોય:
      $M$ એ $\lambd[x][N]$ અને $M'$ એ $\lambd[x][N']$,
      જ્યાં $N \bredpar N'$. બતાવવું છે કે
      $\Subst{(\lambd[x][N])}{R}{y} \bredpar
      \Subst{(\lambd[x][N'])}{R'}{y}$; એટલે કે,
      યોગ્ય તાજા બાંધનાર~$x$ પસંદ કર્યા પછી
      $\lambd[x][\Subst{N}{R}{y}] \bredpar
      \lambd[x][\Subst{N'}{R'}{y}]$.
      આગમનની ધારણા અને \olref{defn:bredpar2}થી આ મળે છે.
      \sourcecorrection{OLLAM-034}{સ્થિર અંગ્રેજી મૂળ
      આ કિસ્સાના બીજા અમૂર્તીકરણમાં
      $\Subst{N'}{R}{y}$ લખે છે, જોકે દાવામાં બીજો
      પ્રતિસ્થાપક~$R'$ છે. અહીં~$R'$ કર્યું છે;
      બાંધનાર તાજો પસંદ કરવાની જરૂર પણ સ્પષ્ટ કરી છે.}
    \item છેલ્લો નિયમ \olref{defn:bredpar3} હોય: અભ્યાસપ્રશ્ન.
    \item છેલ્લો નિયમ \olref{defn:bredpar4} હોય:
      $M$ એ $(\lambd[x][N])Q$ અને
      $M'$ એ $\Subst{N'}{Q'}{x}$. બતાવવું છે કે
      $\Subst{((\lambd[x][N])Q)}{R}{y} \bredpar
      \Subst{\Subst{N'}{Q'}{x}}{R'}{y}$.
      બાંધનાર~$x$ને $y$, $R$ અને~$R'$થી તાજો
      પસંદ કર્યા પછી દાવાને નીચેના સ્વરૂપમાં લાવવાનો રહે છે:
      $(\lambd[x][\Subst{N}{R}{y}])\Subst{Q}{R}{y} \bredpar
      \Subst{\Subst{N'}{R'}{y}}{\Subst{Q'}{R'}{y}}{x}$.
      અહીં જમણી બાજુ માટે સ્થાનાપન્ન-સંયોજનનો
      અલગ લેમ્મો જોઈએ. તે સાબિત થાય તો, આગમનની ધારણા અને
      \olref{defn:bredpar4}થી દાવો મળે છે.
      \sourcecorrection{OLLAM-035}{સ્થિર અંગ્રેજી મૂળ
      જમણી બાજુનાં બે ક્રમિક સ્થાનાપન્નોનો ક્રમ
      કોઈ દલીલ વિના ફેરવે છે. આ માટે~$x\neq y$
      અને~$x$ પ્રતિસ્થાપકમાં મુક્ત ન હોય તેવી
      તાજગી-શરતો સાથે સ્થાનાપન્ન-સંયોજનનો લેમ્મો
      સાબિત કરવો પડે. અહીં તે અનુપસ્થિત સહાયક
      લેમ્મો સ્પષ્ટ કર્યો છે; તેથી સ્થિર મૂળની આ
      સાબિતી આ બિંદુએ અધૂરી છે.}
  \end{enumerate}
\end{proof}

\begin{prob}
  \olref[lam][cr][pb]{lem:comp}ની સાબિતી પૂર્ણ કરો.
\end{prob}

\begin{lem}\ollabel{lem:cont}
  $M \bredpar M'$ હોય તો $M' \bredpar \bcd{M}$.
\end{lem}
\begin{proof}
  $M \bredpar M'$ની !!{derivation} પર આગમનથી.
  \begin{enumerate}
    \item છેલ્લો નિયમ \olref{defn:bredpar1} હોય: અભ્યાસપ્રશ્ન.
    \item છેલ્લો નિયમ \olref{defn:bredpar2} હોય:
      $M$ એ $\lambd[x][N]$ અને $M'$ એ $\lambd[x][N']$,
      જ્યાં $N \bredpar N'$. બતાવવું છે કે
      $\lambd[x][N'] \bredpar
      \bcd{(\lambd[x][N])}$, એટલે
      \olref{defn:bcd2} મુજબ
      $\lambd[x][N'] \bredpar \lambd[x][\bcd{N}]$.
      આગમનની ધારણા અને \olref{defn:bredpar2}થી આ મળે છે.
    \item છેલ્લો નિયમ \olref{defn:bredpar3} હોય:
      $M$ એ $PQ$ અને $M'$ એ $P'Q'$,
      એવા કોઈક $P$, $Q$, $P'$ અને $Q'$ માટે, જ્યાં
      $P \bredpar P'$ તથા $Q \bredpar Q'$.
      આગમનની ધારણાથી $P' \bredpar \bcd{P}$
      અને $Q' \bredpar \bcd{Q}$ છે.
      \begin{enumerate}
        \item $P$ એ $\lambd[x][N]$ હોય તો,
          એવા કોઈક $x$ અને $N$ માટે,
          $P'$ એ $\lambd[x][N']$ એવા કોઈ~$N'$ માટે,
          જ્યાં $N \bredpar N'$. આગમનની ધારણાથી
          $N' \bredpar \bcd{N}$ અને
          $Q' \bredpar \bcd{Q}$. તેથી
          \olref{defn:bredpar4} મુજબ
          $(\lambd[x][N'])Q' \bredpar
          \Subst{\bcd{N}}{\bcd{Q}}{x}$.
        \item $P$ કોઈ $\lambd$-અમૂર્તીકરણ ન હોય તો
          \olref{defn:bredpar3} મુજબ
          $P'Q' \bredpar \bcd{P}\bcd{Q}$;
          \olref{defn:bcd3} મુજબ જમણી બાજુ
          $\bcd{PQ}$ છે.
      \end{enumerate}
    \item છેલ્લો નિયમ \olref{defn:bredpar4} હોય:
      $M$ એ $(\lambd[x][N])Q$ અને
      $M'$ એ $\Subst{N'}{Q'}{x}$,
      એવા કોઈક $x$, $N$, $Q$, $N'$ અને~$Q'$ માટે, જ્યાં
      $N \bredpar N'$ અને $Q \bredpar Q'$.
      આગમનની ધારણાથી $N' \bredpar \bcd{N}$
      અને $Q' \bredpar \bcd{Q}$. તેથી
      \olref{lem:comp} મુજબ
      $\Subst{N'}{Q'}{x} \bredpar
      \Subst{\bcd{N}}{\bcd{Q}}{x}$;
      જમણી બાજુ બરાબર
      $\bcd{((\lambd[x][N])Q)}$ છે. અહીં
      \olref{lem:comp}ની અગાઉ નોંધેલી સાબિતી-ખામીનો
      આધાર હજી ખુલ્લો છે.
  \end{enumerate}
\end{proof}

\begin{prob}
  \olref[lam][cr][pb]{lem:cont}ની સાબિતી પૂર્ણ કરો.
\end{prob}

\begin{thm}\ollabel{thm:cr}
  $\bredpar$ ચર્ચ--રોસર ગુણધર્મ ધરાવે છે.
\end{thm}

\begin{proof}
  \olref{lem:cont}માંથી તરત મળે છે; એ લેમ્માની
  ઉપર નોંધેલી સાબિતી-ખામી અહીં પણ લાગુ પડે છે.
\end{proof}

\end{document}
